\begin{proposition}[Properness criterion]
\label{prop:properness-biholomorphism}
Let \(I:A_s\to A_t\) be a proper holomorphic immersion between
connected annuli.  If its oriented annular degree is one, then \(I\)
is a biholomorphism.  In particular, every universal-cover lift of
\(I\) is injective.
\end{proposition}

\begin{proof}
A holomorphic immersion between Riemann surfaces is a local
biholomorphism; therefore \(I(A_s)\) is open in \(A_t\).  A proper
continuous map from a locally compact Hausdorff space to a Hausdorff
space is closed; hence \(I(A_s)\) is also closed in \(A_t\).  The
target annulus is connected and the image is nonempty, so
\[
 I(A_s)=A_t .
\]

A proper local homeomorphism is a covering map.  For completeness:
each target point has a discrete fiber by local injectivity and a
compact fiber by properness, hence a finite fiber.  Choose pairwise
disjoint neighborhoods of the finitely many preimages on which \(I\)
is invertible; after the target neighborhood is shrunk, properness
excludes additional preimages, so the target point has an evenly
covered neighborhood.

The number \(N\) of sheets of a connected covering of annuli is the
index of the induced subgroup
\[
 I_*\pi_1(A_s)\subset\pi_1(A_t).
\]
For an orientation-preserving holomorphic annulus covering, \(I_*\) is
multiplication by \(N\).  The assumed degree is one, so \(N=1\), and
\(I\) is bijective.  A bijective local biholomorphism has a
holomorphic inverse, so \(I\) is a biholomorphism.  The final
assertion follows from Theorem~\ref{thm:tile-base-equivalence}.
\end{proof}